\section{Logic programs as argumentation: complete proofs}
\label{sec:annex-p18-logic-program-argumentation}

\begin{annexitem}{annex:p18.support.reduct}{thm:p18-support-reduct}{Finite proof witnesses}
\begin{proof}
The positive consequence operator is monotone, and its iterates
from the empty set increase. A finite body contained in their
union is already contained in one finite stage: take the maximum
of the finitely many premise stages. Its head belongs to the
next stage. An empty body gives its head at stage one. The
union is therefore fixed. Every model, meaning a set closed
under the clauses, contains every stage by induction, proving
leastness. Induction on the stage constructs a finite proof
tree; induction on tree height puts each finite proof's root
in a finite stage. Finite derivation sequences and trees
convert into each other by unfolding earlier derivations
or reading a tree in postorder.

A surviving reduct clause has all its negative premise atoms
outside \(X\). Restoring those leaves in a finite reduct proof
gives a finite support in \(\neg(H\setminus X)\). Conversely,
deleting the negative leaves of such a support proof gives
a reduct proof. This proves the interface with infinitely
many ground clauses as well. If \(X\subseteq Y\), the
reduct for \(Y\) has no clause absent from the reduct
for \(X\), so least-model containment gives
\(\Gamma(Y)\subseteq\Gamma(X)\).
\end{proof}
\end{annexitem}

\begin{annexitem}{annex:p18.support.operator}{thm:p18-support-operator}{The literal support-graph identities}
\begin{proof}
All supports here are negative, so negative unit arguments
attack nothing. An argument with support \(K\) is attacked
by \(E\) exactly when \(\neg p\in K\) for some \(p\in T\).
This proves \(N_\infty(E)={\cal E}_\infty(H\setminus T)\).
Its positive conclusions are \(\Gamma(T)\) by the reduct
interface, and its negative conclusions are
\(\neg(H\setminus T)\). The characteristic operator of
any framework is the no-attack operator squared, by
\cref{thm:p18-alternating-levels}. Applying the calculation
twice proves both \(F_\infty\) formulas. All finite
supports, including unused assumptions, were retained.
If \(p,\neg p\) were conclusions of a conflict-free
set, an included positive argument for \(p\) would
attack the included unit \(\nu_p\). Hence conclusions
are consistent. For admissible \(E\), taking conclusions
in \(E\subseteq F_\infty(E)\) gives the two remaining
inclusions.
\end{proof}
\end{annexitem}

\begin{annexitem}{annex:p18.program.stable.correspondence}{thm:p18-program-stable-correspondence}{Both stable-model directions}
\begin{proof}
If \(M=\Gamma(M)\), put
\(E={\cal E}_\infty(H\setminus M)\).
Its conclusions are \(\widehat M\), and the operator
identity gives \(N_\infty(E)=E\). The no-attack
characterization therefore makes it stable.
Conversely, for a stable \(E\), let \(M\) be its
positive conclusions. Stability and the identity
force \(E={\cal E}_\infty(H\setminus M)\); taking
positive conclusions gives \(M=\Gamma(M)\).
The negative conclusions are exactly the complementary
ones. These calculations also give the two inverse
maps and uniqueness of each corresponding extension.
\end{proof}
\end{annexitem}

\begin{annexitem}{annex:p18.unfounded.support}{thm:p18-unfounded-support}{The full unfounded-set and Appendix B argument}
\begin{proof}
A union of unfounded sets is unfounded: for its head
choose one constituent containing that head. Every
clause is false or has a positive premise in that
constituent, hence in the union. This proves greatestness.
Enlarging \(I\) preserves falsity witnesses, so makes
\(U_Q\) monotone. It also preserves truth witnesses,
giving monotonicity of \(T_Q,V_Q,T_Q^*,W_Q\).
These statements remain valid on the auxiliary full
literal powerset.

Let \(E\) be admissible, \(I=\operatorname{con}(E)\),
\(T=I^+\), \(D=\Gamma(T)\), and \(X=H\setminus D\).
The support-operator lemma gives \(I^-\subseteq X\).
The set \(X\) is unfounded. A clause with head
outside \(D\) is either deleted in \(Q^T\), in
which case a negative premise is false in \(I\),
or survives. A surviving clause must have a
positive premise outside \(D\), since \(D\) is
closed under that clause. Thus \(X\subseteq U_Q(I)\)
and \(I^-\subseteq U_Q(I)\).

We prove the full finite-tree claim of Appendix B:
a proof tree whose root is in \(U_Q(I)\) has a
path to a negative leaf \(\neg t\) with \(t\in T\),
all its positive nodes lying in \(U_Q(I)\).
At a current positive node inspect its clause.
If a negative premise is false in \(I\), choose
that leaf and stop. If a positive premise \(b\)
is false, then \(b\in I^-\subseteq U_Q(I)\);
continue down its subtree. If the body is not
false, unfoundedness gives a positive child in
\(U_Q(I)\), and we continue there. A fact cannot
expand a node in an unfounded set, since it
has neither a false body nor a positive child.
Each continuation strictly decreases remaining
finite tree height, so it terminates at the
specified negative leaf. This proves all
negative-false, positive-false, unfounded-child
and fact cases.

If a member of \(U_Q(I)\) belonged to \(D\),
the reduct interface would supply a tree with
every negative leaf's atom outside \(T\), contrary
to the tree claim. Therefore \(U_Q(I)\subseteq X\),
proving equality and, by the operator formula,
the unit-argument characterization of Lemma B.1.
The selected leaf also gives an actual attack
on every proposed support proof of an unfounded
root, retaining the source's proof-tree mechanism.

Put \(B=U_Q(I)=H\setminus\Gamma(T)\).
Positive saturation from \(\neg B\) gives
\(\Gamma(H\setminus B)=\Gamma^2(T)\).
Admissibility gives \(T\subseteq\Gamma^2(T)\)
and \(I^-\subseteq B\). Hence adding the seed
\(I\) does not change that closure: it contains
the closure from \(\neg B\), while
\(\Gamma^2(T)\cup\neg B\) contains the seed
and is closed under positive consequence.
These two inclusions prove
\[
 W_Q(I)=\Gamma^2(T)\cup\neg B
                    =\operatorname{con}(F_\infty(E)).
\]
Finite bodies justify the countable saturation.
\end{proof}
\end{annexitem}

\begin{annexitem}{annex:p18.program.well.founded}{thm:p18-program-well-founded}{The two least fixed points and the ordinal correspondence}
\begin{proof}
For a monotone \(f:2^L\to2^L\), intersect all pre-fixed
sets \(J\), meaning \(f(J)\subseteq J\), and call
the intersection \(K\). The family is nonempty
because it contains \(L\). Monotonicity gives
\(f(K)\subseteq K\), and then
\(f(f(K))\subseteq f(K)\). Thus \(f(K)\)
is in the intersected family, so \(K\subseteq f(K)\).
This proves a least fixed point, below every
pre-fixed point. Its ordinal iteration from
the empty set, with limit unions, stabilizes
by the distinct-new-element cardinality argument
of \cref{annex:p18.transfinite.grounded}, and
is this least fixed point.

Apply this to \(V_Q,W_Q\). Since \(W_Q\)
contains its input, \(W_Q(I)=I\) exactly
when \(V_Q(I)\subseteq I\): the latter
condition supplies the new negative assumptions
and closes \(I\) under \(T_Q\).
Every \(V_Q\)-fixed point is therefore
\(W_Q\)-fixed. Conversely the least
\(W_Q\)-fixed point is pre-fixed for
\(V_Q\). Leastness in the two directions
proves equality of their least fixed points.

Every grounded argument stage \(E_\alpha\)
is admissible. The preceding lemma supplies
the exact successor identity
\(\operatorname{con}(E_{\alpha+1})
=W_Q(\operatorname{con}(E_\alpha))\).
Conclusions commute with unions at limits:
an included literal has an actual argument
in some earlier stage. The empty stages
agree, so transfinite induction proves
the stated stage correspondence. At
stabilization it gives the common least
fixed point. Grounded conflict-freeness
makes its conclusions consistent, proving
the promised partial-interpretation property.
\end{proof}
\end{annexitem}

\begin{annexitem}{annex:p18.clause.completion}{thm:p18-clause-completion}{The exact completion and extension isomorphism}
\begin{proof}
For a ground head, the completed body's existential
substitutions are exactly the ground instances of
the source clauses, by identity of Herbrand terms.
Its disjunction is true when some body is true,
false when every body is false, and undefined
otherwise. The empty disjunction is false.
Thus equality of all three truth values is
precisely \(I=\Phi_Q(I)\); in the total case
it is the ordinary supported-model equation.
Free Herbrand terms satisfy the stated equality
axioms. No infinite ground disjunction has
been treated as a finitary formula.

For any \(E\), let \(I=\operatorname{con}(E)\)
and \(Z=Z_Q(I)\). Every positive argument
for \(p\) is attacked exactly when each
\(p\)-clause has a false body. Therefore
\[
 \nu_p\in F_{\rm c}(E)\quad\Longleftrightarrow\quad p\in Z.
\]
A clause argument with body \(K\) is defended
exactly when each positive premise \(q\)
has its unit attacker \(\nu_q\) attacked,
meaning \(q\in I^+\), and every positive
argument for each negative premise atom
\(q\) is attacked, meaning \(q\in Z\).
These are all possible attackers under
the literal clause-graph definition.

If \(E\) is complete, \(E=F_{\rm c}(E)\).
Its conclusions are consistent, since
both signs of \(p\) would give an
attack on its included unit. The unit
calculation gives \(I^-=Z_Q(I)\);
the clause calculation then says the
included clause arguments are exactly
those with bodies true in \(I\).
Thus \(I^+=T_Q(I)\), \(I=\Phi_Q(I)\),
and \(E={\cal E}_{\rm c}(I)\).
Conversely let \(I\) be a consistent
\(\Phi_Q\)-fixed point. The conclusions
of \({\cal E}_{\rm c}(I)\) are exactly
\(I\). An attack within that set would
put an included conclusion and its
complement from an included support
in \(I\), contradicting consistency.
The same defense calculation gives
\(F_{\rm c}({\cal E}_{\rm c}(I))
={\cal E}_{\rm c}(I)\), so it is
complete. Truth of bodies and taking
conclusions are monotone; the inverse
maps therefore preserve and reflect order.

Preferred extensions are exactly maximal
complete extensions. A preferred set is
complete; if a maximal complete set had
a proper admissible superset, that
superset would extend to a larger
preferred, hence complete, set. This
proves the converse and transfers
maximality through the isomorphism.

A stable extension has total conclusions:
if it has no positive conclusion \(p\),
the unit \(\nu_p\) has no attacker
in it and must be included by stability.
Conflict-freeness excludes both signs.
Conversely for total \(I\), every
omitted clause argument has a false
body, hence an included complementary
conclusion attacking it. Every omitted
unit has its atom true and therefore
an included positive argument attacking
it. This proves stability and the
two directions of Theorem 55.
\end{proof}
\end{annexitem}

\begin{annexitem}{annex:p18.clause.fitting}{thm:p18-clause-fitting}{Correct successor stages and the Fitting fixed point}
\begin{proof}
The preceding object-level defense calculation applies
to every \(E\). Its unit conclusions are
\(\neg Z_Q(I)\), and its positive conclusions
are exactly the heads with positive premises
in \(I^+\) and negative premise atoms in
\(Z_Q(I)\). This is the stated \(\Psi_Q(I)\).
Falsity witnesses are preserved by information
inclusion, so \(Z_Q,\Phi_Q,\Psi_Q\) are monotone
on the full literal powerset.

If \(I=\Phi_Q(I)\), then \(I^-=Z_Q(I)\)
and \(I^+=T_Q(I)\), so \(\Psi_Q(I)=I\).
Conversely \(I=\Psi_Q(I)\) forces the
same negative equation, after which its
positive equation becomes \(I^+=T_Q(I)\).
Thus the two operators have exactly
the same fixed points and the same
least fixed point. Successor stages
of the argument iteration equal
\(\Psi_Q\) stages by the proved identity,
and limit conclusions commute with
union. Their least fixed point is
therefore the grounded conclusion set,
which is consistent. This proves the
Fitting correspondence without asserting
equal successor stages for \(\Phi_Q\).
\end{proof}
\end{annexitem}

\begin{annexitem}{annex:p18.used.support.reduction}{thm:p18-used-support-reduction}{The complete bridge retaining the literal support graph}
\begin{proof}
Every finite support proof uses an actual finite set \(J\)
of negative leaves contained in its declared support \(K\).
Thus the used graph has exactly the same derivability
interface with \(\Gamma\): a reduct proof exists if
and only if a used-support proof with permitted leaves
exists. The calculations of \(N_\infty,F_\infty\)
therefore apply verbatim to the used graph, with
its own argument universe.

In either graph let \(E\) be complete and let
\(B=\operatorname{con}(E)^-\). A positive argument
with support \(K\) is defended exactly when each
unit \(\nu_q\) for \(\neg q\in K\) is defended,
since both conditions quantify over the identical
positive arguments for \(q\). Completeness makes
this equivalent to \(K\subseteq\neg B\).
The positive conclusions are consequently
\(T=\Gamma(H\setminus B)\).
The unit defense calculation says
\(B=H\setminus\Gamma(T)\).
Conflict-freeness gives \(B\cap T=\varnothing\).
These are the two displayed conditions.

Conversely, if \(B\) satisfies both conditions,
take every argument of the chosen universe whose
support is contained in \(\neg B\), including
the specified units. Its conclusions are
\(T\cup\neg B\) and are consistent.
Any attack within the set would require a
positive conclusion in \(T\) whose negative
literal belongs to \(\neg B\), contradicting
consistency. The unit and positive defense
calculations then make the set fixed, hence
complete. This proves the exact parametrization.

Inclusion of two such complete sets is equivalent
to inclusion of their \(B\)'s, because each
\(B\)-membership is recorded by a unit argument.
Both universes therefore give the same ordered
parameter family and the same conclusions.
Least complete and maximal complete elements
correspond, so grounded and preferred extensions
correspond. Stability is exactly totality
\(B=H\setminus T\), by the stable-model
calculation in either graph, and is preserved.
Finally the intersection of any such argument
family consists exactly of the arguments whose
supports lie in the intersection of its \(B\)'s.
Thus equality of the grounded set with the
preferred intersection holds in one graph
exactly when it holds in the other.
\end{proof}
\end{annexitem}

\begin{annexitem}{annex:p18.program.graph.transfer}{thm:p18-program-graph-transfer}{All signed-walk and rank cases}
\begin{proof}
For a used-support argument with positive conclusion
predicate \(p\), every negative leaf \(\neg q(\mathbf t)\)
has a path from the root through positive premises
to that leaf. It gives a predicate dependency walk
from \(p\) to \(q\) with exactly one negative edge.
This assertion concerns actual leaves; it is why
unused assumptions cannot be included in this step.

Assign sign \(0\) to a positive argument conclusion
and sign \(1\) to a negative unit conclusion.
For an attack from \(B\) to \(A\), the predicate
walk goes from the predicate of \(A\) to that
of \(B\). If both conclusions are positive,
it has odd negative parity: use the previous
root-to-leaf walk in the used graph, or the
negative body edge in the clause graph.
If \(B\) is a negative unit attacking a
positive clause argument, use its positive
body edge, of even negative parity.
If \(A\) is a negative unit, \(B\) is
positive with the same predicate and we
use the zero-length walk. There is no
negative-to-negative attack. Thus in
both graphs an attack walk of length \(m\)
from \(B\) to \(A\) gives a dependency
walk with parity
\[
                         m+\operatorname{sign}(A)+\operatorname{sign}(B)\pmod2.
\]
This follows by concatenation, since the
signs at internal vertices cancel.

If \(B\) is controversial with respect
to \(A\), two opposite-parity attack
walks give opposite-parity dependency
walks with the same endpoint predicates.
If both are nonempty, strictness excludes
them. If one is empty, the predicates
coincide and the other is an odd-negative
closed walk. Traversing it twice gives
a nonempty even-negative closed walk,
again violating strictness. This proves
the two uncontroversial conclusions.

For call-consistency, the same argument
shows that a controversial pair either
already gives an odd-negative closed
walk, or gives nonempty dependency walks
of both parities between distinct
predicates. An infinite backward
controversy sequence would repeat a
predicate, since the source program
has finitely many predicates. Between
two repetitions, choose an odd-parity
walk for one link and even-parity
walks for every other link. Their
concatenation is an odd-negative
closed dependency walk, a contradiction.
Thus both graphs are limited controversial.

For stratification, collapse the finite
predicate graph into strongly connected
components. No negative edge remains
inside a component. The finite acyclic
component graph has integer ranks
nonincreasing along all dependency
edges and strictly decreasing across
components. Every root-to-negative-leaf
walk of a used support crosses a
negative edge, and hence strictly
decreases rank. In a backward attack
sequence in the used graph, all
arguments after the first are positive:
negative units attack nothing there.
Every subsequent step thus strictly
decreases a nonnegative integer rank.
No infinite sequence exists.

For a hierarchical program the entire
predicate graph is acyclic, so choose
ranks strictly decreasing on every
dependency edge. In the clause graph,
every backward attack step strictly
decreases rank except the step from
a negative unit to a positive attacker
with the same predicate. Two such
zero-length steps cannot be consecutive,
because the intermediate argument is
positive. An infinite backward sequence
would still contain infinitely many
strict decreases, which is impossible.

Finally the fact \(p\leftarrow\) is
stratified, hierarchical, strict and
call-consistent. In the literal support
graph, \(\{\neg p\}\) supports \(p\)
by ignoring that assumption. Its
argument attacks itself by the exact
attack clause. Repeating the self-attack
violates well-foundedness and limited
controversy, and the odd self-walk
together with the zero even walk
violates uncontroversiality. This proves
the stated boundary of the graph
assertions and the necessity of the
explicit semantic reduction.
\end{proof}
\end{annexitem}

\begin{annexitem}{annex:p18.program.coincidence}{thm:p18-program-coincidence}{Every coincidence and existence consequence}
\begin{proof}
For a stratified program the used graph is
well-founded, so its unique complete extension
is grounded and stable. The reduction transfers
these facts to the literal support graph.
The stable and well-founded correspondences
therefore give the same unique total program
interpretation. Hierarchy applies the same
well-founded theorem directly to the clause
graph; its completion and Fitting
correspondences prove the second assertion.

For strictness, the used and clause graphs
are uncontroversial, so they are coherent
and relatively grounded. The reduction
transfers coherence and relative grounding
to the literal support graph.
For the following residual argument work
in the uncontroversial used graph; its
grounded conclusions and stable-model
correspondence are preserved by the reduction.
To prove the literal query assertion,
one must not exchange an intersection
with existential choice of a proof.
Instead use the unique unit \(\nu_p\).
If \(p\) is not a grounded conclusion
and \(\neg p\) is grounded, every stable
conclusion set makes \(p\) false.
If neither sign is grounded, \(\nu_p\)
lies in the residual: it is absent
and no grounded positive argument
attacks it. The two residual-extension
alternatives of
\cref{thm:p18-relative-grounding}
give a complete extension containing
\(\nu_p\), and one attacking \(\nu_p\).
Extend each to a preferred extension;
coherence makes each stable. Their
total conclusions make \(p\) respectively
false and true. Finally every grounded
literal belongs to every stable
extension's conclusion set by containment.
These cases prove both positive and
negative skeptical assertions. Preferred
and stable extensions coincide by coherence.

Call-consistency gives limited controversy
for the used and clause graphs, hence
coherence and stable existence there.
Transfer through the reduction gives
coherence of the literal support graph.
The clause-graph isomorphism sends
maximal partial completion models to
preferred extensions, now stable, so
their interpretations are total.
Stable existence in the support graph
gives a stable model; stable existence
in the clause graph gives a total
Herbrand completion model. These prove
all listed consequences, including
nonempty model families in the
skeptical statement.
\end{proof}
\end{annexitem}

\begin{annexitem}{annex:p18.program.computability}{thm:p18-program-computability}{Exact finite procedures and a halting reduction}
\begin{proof}
For the clause graph, unit arguments are recognized
directly. For a positive supplied ground pair
\((K,p)\), try each of the finitely many source
clauses. Match its head to \(p\), and assign
each of its finitely many body occurrences to
an element of \(K\), requiring the assigned
image to be exactly \(K\). Check the resulting
finite equations between source terms and
ground terms for a consistent substitution.
There are finitely many assignments and
finite term comparisons. This recognizes
exactly ground instances, even when repeated
body literals collapse in the body set.
It treats facts and empty \(K\) directly.
Thus clause-argument membership is decidable.

For support arguments, enumerate finite ground
derivation sequences over the effectively
enumerable Herbrand terms of a finite source
program. Ground-clause membership and the
finite support checks are decidable, so
this semidecides existence of a proof.
It is not a total decision procedure.
Here is an explicit reduction establishing
the limitation without assuming that every
non-total partial function has an
undecidable graph.

Encode a Turing configuration as
\(\mathrm{cfg}(q,L,a,R)\), with state
constant \(q\), current symbol \(a\),
and finite lists \(L,R\), nearest
neighbor first; use list constructors
\(\mathrm{nil},\mathrm{cons}\) and
blank symbol \(\square\).
Use predicate \(\mathrm{reach}(w,C)\),
retaining the original input word \(w\).
The initial clause schemas are the facts
\[
\begin{gathered}
 \mathrm{reach}(\mathrm{nil},\mathrm{cfg}(q_0,\mathrm{nil},\square,\mathrm{nil}))
                         \leftarrow,\\
 \mathrm{reach}(\mathrm{cons}(a,R),
                    \mathrm{cfg}(q_0,\mathrm{nil},a,R))\leftarrow.
\end{gathered}
\]
For a transition writing \(b\), moving right
and entering \(q'\), include
\[
\begin{aligned}
 \mathrm{reach}(w,\mathrm{cfg}(q',\mathrm{cons}(b,L),c,R))
   &\leftarrow\mathrm{reach}(w,\mathrm{cfg}(q,L,a,\mathrm{cons}(c,R))),\\
 \mathrm{reach}(w,\mathrm{cfg}(q',\mathrm{cons}(b,L),\square,\mathrm{nil}))
   &\leftarrow\mathrm{reach}(w,\mathrm{cfg}(q,L,a,\mathrm{nil})).
\end{aligned}
\]
For a left move use
\[
\begin{aligned}
 \mathrm{reach}(w,\mathrm{cfg}(q',L,c,\mathrm{cons}(b,R)))
   &\leftarrow\mathrm{reach}(w,\mathrm{cfg}(q,\mathrm{cons}(c,L),a,R)),\\
 \mathrm{reach}(w,\mathrm{cfg}(q',\mathrm{nil},\square,\mathrm{cons}(b,R)))
   &\leftarrow\mathrm{reach}(w,\mathrm{cfg}(q,\mathrm{nil},a,R)).
\end{aligned}
\]
A stationary move only changes \(q,a\) to \(q',b\).
For each halting state add
\(\mathrm{halt}(w)\leftarrow
  \mathrm{reach}(w,\mathrm{cfg}(q_h,L,a,R))\).
This is a finite definite program for any
finite transition table. Induction on
derivation length proves that, for a
valid ground input list \(w\), its
derived configurations are exactly the
machine's finitely reached configurations.
Thus \(\mathrm{halt}(w)\) has a finite
proof exactly when the machine halts.

Take a fixed universal machine, which
simulates the finite transition table
encoded in its input. Its halting
language is undecidable: a supposed
halting decider applied to a program
on its own encoding permits a program
that loops exactly when the decider
says it halts, and halts otherwise.
Applying this program to its own
encoding is a contradiction.
The fixed universal-machine clause
program above therefore has an
undecidable set of queries
\((\varnothing,\mathrm{halt}(w))\)
in its support graph. This proves
even a fixed-program instance of
the claimed argument undecidability.

For a definite program let \(D\)
be its least positive model.
Its reduct is the same program
for every interpretation, so
\(\Gamma(X)=D\) for every \(X\).
The support-operator formula makes
the first grounded conclusion set
\(D\cup\neg(H\setminus D)\);
it is total and stable. Grounded
containment and stability imply
that it is the unique complete
and preferred conclusion set.
Thus the same halting query reduces
directly to well-founded, stable,
credulous preferred and skeptical
preferred truth. Argument
undecidability alone was not used
as an implication about semantics.

If the explicitly grounded input
and \(H\) are finite, enumerate
all finite negative subsets \(K\)
and all heads \(p\), and compute
finite positive closure to test
support. This constructs the
entire support graph; the clause
graph is read directly from the
grounded clauses and units.
Their finite subset procedures
then give every extension family.
These are precisely the finite
domains claimed in the statement.
\end{proof}
\end{annexitem}

\begin{annexitem}{annex:p18.program.examples}{thm:p18-program-examples}{All small and Appendix A calculations}
\begin{proof}
For \(Q_{\rm delay}\), no clause has head \(q\),
so the first \(\Phi\) iterate includes \(\neg q\).
The only \(p\)-body is initially undefined, so
it includes neither sign of \(p\). The next
iterate adds \(p\) and is fixed.
In the clause graph the \(p\)-clause argument
and \(\nu_q\) have no attackers; both belong
to its first grounded stage. They attack
the remaining unit \(\nu_p\). This stage
is fixed and gives \(p,\neg q\).

For \(Q_{\rm loop}\) there is no finite
positive proof of \(p\), so the support
graph has only the unit \(\nu_p\).
It is grounded and stable, giving
the well-founded value \(\neg p\).
The positive reduct has least model
empty, hence that is its sole stable
program model. Completion is \(p=p\)
as a truth-value equation, so its
total models are \(p\) false and
\(p\) true, and its least partial
model leaves \(p\) undefined.
Equivalently its clause graph is
the mutual attack pair consisting
of \((\{p\},p)\) and \(\nu_p\).
This also gives the stated Fitting value.
For the fact program, \((\varnothing,p)\)
is unattacked and defeats both the
unused-support argument \((\{\neg p\},p)\)
and \(\nu_p\); its sole complete
and stable conclusions are \(\{p\}\).
The clause graph has the fact and
the defeated unit, with the same
conclusions.

In Appendix A, each distinct
\((\{\neg q_n\},p)\) is both a
used support argument and a
ground-clause argument. All attack
\(\nu_p\), so both source graphs
have an argument with infinitely
many attackers and are non-finitary.
For the displayed support-graph
stage formula, at stage zero
\(\Gamma(\varnothing)=H\), because
every negative premise is permitted;
the second reduct then derives only
\(e_0,q_0\). This gives \(I_1\).
For the positive part \(T_m=I_m^+\),
\(m\geq1\), direct reduct evaluation gives
\[
 \Gamma(T_m)=
 \{q_n:n<\omega\}\cup\{p,r\}
 \cup\bigl(\{e_n:n<\omega\}\setminus
                 \{e_{2j+1}:j<m\}\bigr).
\]
Indeed an \(e_n\) clause is removed
exactly when its predecessor even
atom is already in \(T_m\).
Every \(q_n\) follows either from
its surviving negative clause or
from its positive \(e_n\) clause.
Some \(q_n\) is beyond the finite
stage, allowing \(p\); and \(p\)
is absent from \(T_m\), allowing
\(r\). Reducing once more removes
all \(p\)-clauses and the \(r\)-clause,
and derives exactly the even \(e\)'s
through index \(2m\) and the
\(q\)'s through index \(2m\).
The complementary negative part
is exactly the odd \(e\)'s
through index \(2m-1\).
The support-operator identity
therefore gives the formula for
\(I_{m+1}\), completing induction.

At the union stage all \(q_n\)
are true and the \(e_n\)'s
have their alternating signs.
The reduct of its positive
part contains no surviving
\(p\)-clause and does contain
\(r\leftarrow\); its least
model is all \(q_n\), all
even \(e_n\), and \(r\).
A second reduction gives
the same positive model.
The operator identity therefore
adds exactly \(\neg p,r\)
at the next stage. That
interpretation is total and
fixed. Neither literal was
in any finite stage or
their union, proving the
exact \(\omega+1\) boundary.
\end{proof}
\end{annexitem}
